\thispagestyle{plain}

\begin{center}
    \Large
    \thesistitle

    \vspace{0.4cm}
    \large
    \thesisauthorname
       
    \vspace{0.4cm}
    \large
    M.Sc. Research Proposal
      
    \vspace{0.2cm}
    \large
    Ben-Gurion University of the Negev 
    
    \vspace{0.2cm}
    \large
    \thesisyear
    
    \vspace{0.9cm}
    \Large
    \textbf{Abstract}
\end{center}

The Model Context Protocol (MCP) has become the standard interface through
which autonomous large-language-model agents invoke external tools ---
filesystems, databases, code hosts, and messaging platforms. Almost all
existing MCP-security research protects the \emph{agent} from a malicious
server. This research proposal inverts that direction and takes the standpoint
of the \textbf{server as the protected asset and the agent as the threat
source}: a server that fronts sensitive resources has no native way to decide,
when a call arrives, whether it should be allowed, throttled, or denied, and an
indirectly-injected agent is indistinguishable from a legitimate one.

We propose \textbf{McpRisk}, a defense-oriented
risk-scoring framework that assigns every MCP tool interaction a graduated,
interpretable four-band risk score. The score is a product of three
interpretable factors --- the sensitivity of the asset, the blast radius of the
operation, and the impact (reversibility) of the tool, the last grounded in a
13-operation atomic-operation severity taxonomy --- computed statically at
design time and escalated at request time by a parameter layer that reads the
actual arguments of a call. The pipeline is driven by a single deterministic
local model and audited by an independent review subsystem.

A proof-of-concept implementation and preliminary evaluation against four kinds
of ground truth --- realistic MCP deployments, a ten-rater human-and-framework
oracle panel, a corpus of captured agent calls, and external attack benchmarks
--- already show that the approach is feasible: McpRisk agrees with reviewed
judgments on tool impact $97\%$ of the time and attains a quadratic-weighted
$\kappa$ of $+0.66$ against the human panel, and, unlike capability-only scorers,
retains discriminative power on live external attacks. The preliminary results
also expose the challenges the remaining research must solve, chiefly the
fragility of the composite risk band. This proposal sets out the scoring model,
these preliminary results, and a detailed research plan for developing a
contextual request-time scorer and a deployable server-side risk gate.



